\section{Podsumowanie}
W ramach badania przygotowano model uczonej kompresji stratnej obrazów, adaptując do tego zadania architekturę transformera wizyjnego z oknem przesuwnym. Do realizacji sieci neuronowej wykorzystano otwartoźródłowy framework PyTorch, oferujący narzędzia do tworzenia rozwiązań opartych na głębokim uczeniu maszynowym. Wykorzystano mieszany zbiór treningowy zawierający obrazy o wysokiej rozdzielczości, przechowywanych w formacie bezstratnym, aby uniknąć modelowania statystyk artefaktów kompresji stratnej. Przeprowadzono kilkadziesiąt eksperymentów w celu znalezienia optymalnej architektury, hiperparametrów oraz technik uczenia maszynowego, doprowadzając do uzyskania finalnego modelu. Otrzymano kodek zdolny do konkurowania z ugruntowanymi klasycznymi metodami kompresji.
\subsection{Wnioski z badania}
\subsection{Możliwe usprawnienia}